\chapter{Pressure projection}\label{ch:projection}
Thesis Chapter~7 presents one- and two-velocity pressure operators,
CMDGFM, and a mimetic gradient~\cite{remmerswaal2023}. This chapter will
document the actual projection systems and boundary conditions.

\section{One-velocity projection}
The projection is non-incremental: it solves for the pressure rather than a
pressure increment.
\begin{implnote}
\apitype{ProjectionWorkspace}, \apifun{project!}, \apitype{OneVelocityModel}.
\end{implnote}

\section{CMDGFM jump interpolation}
The unknown $\xi$ is the scaled pressure-gradient jump; $\beta$ controls its
interpolation. Write $R=|n_t/n_f|$ for the ratio of tangential to normal
interface-normal components, $s$ for the sum of the cross-phase inference-face
mass fractions, and $x=|\stagger{\+\eta}\cdot\+n_f|$ for the normal alignment.
\begin{deviation}{CMDGFM beta}
The code uses $\beta=(1-s)/(R-s)$ for $R>1$ (with clamping), rather than the thesis's $3x^2-2x^3$.
\end{deviation}
\begin{deviation}{CMDGFM inference and unresolved rows}
Inference faces use extreme mass fractions among same-axis vertex neighbours; unresolved jump rows set $\xi=0$.
\end{deviation}
\begin{deviation}{CMDGFM linear solver}
The implementation defaults to BiCGSTAB with a tree-GMG/active-jump-ILU block preconditioner, which is not the thesis's solver.
\end{deviation}
\begin{implnote}
\apitype{CMDGFMProjection}, \apitype{CMDGFMProjectionSolver},
\apifun{cmdgfm_jump_interpolant}, \apifun{cmdgfm_gradient!}.
\end{implnote}

\section{Mimetic projection}
The implemented local constraints and active-face rule will be described with
their energy and momentum identities.
\begin{deviation}{Mimetic KKT system}
The code solves a locally constrained $(p,\xi,\lambda)$ KKT system; the thesis solves for potentials $\psi_\eta$ and $(p,\psi_\tau)$.
\end{deviation}
\begin{implnote}
\apitype{MimeticProjection}, \apitype{MimeticProjectionSolver},
\apifun{mimetic_gradient!}.
\end{implnote}

\section{Velocity and pressure boundaries}
Inflow and outflow boundary handling, the pressure condition, and the
impermeability requirement of WY transport will be stated separately.
\begin{implnote}
\apitype{VelocityBoundaryConditions}, \apitype{Inflow}, \apitype{Outflow},
\apitype{ZeroGradientVelocity}.
\end{implnote}
